\title{Urban Fabric and Climate Adaptation Priorities in \.{I}zmir}

\author[a]{Yusuf Emino\u{g}lu}
\author[b]{Halil Top\c{c}u}
\author[a]{Hilmi Evren Erdin}
\affil[a]{Department of City and Regional Planning, Dokuz Eyl\"ul University,
\.{I}zmir, T\"urkiye}
\affil[b]{Graduate School of Natural and Applied Sciences, Urban Design Program,
\.{I}zmir Demokrasi University, \.{I}zmir, T\"urkiye}

\date{\today}

\maketitle

\begin{center}
\small Corresponding author: yusuf.eminoglu@deu.edu.tr
\end{center}

\begin{abstract}
Urban form shapes climate hazard, yet reproducible morphometrics and resilience
assessment rarely meet at the neighbourhood scale of adaptation. This study
couples street-scale morphometrics with multi-hazard
exposure and age-structure social vulnerability to identify fabric-specific
adaptation priorities for the \.{I}zmir functional urban region. An open-source,
parameter-logged QGIS workflow analyzes a 250\,m grid covering all 3{,}777 urban
cells, stratified into seven a-priori fabric types. A catchment-radius
cross-attribution rule combines area-weighted tessellation characters with
400/800\,m network service-area reaches in a common unit. An explainable
gradient-boosting model with SHAP relates measured land-surface temperature to
morphological mechanisms. Pareto optimization, TOPSIS and entropy/Monte-Carlo
weight testing then rank priorities across five need axes: heat, cooling deficit,
access deficit, coastal exposure and social vulnerability. Pluvial susceptibility
remains a screening overlay pending hydrodynamic validation, while Moran's $I$,
LISA, Getis-Ord $G_i^{*}$ and a Gini measure describe spatial structure and equity.
Across the 3{,}777 cells, the seven fabrics consolidate into four morphometric
super-types whose broad structure persists across grid resolution, although the
finest mode is scale-sensitive (Adjusted Rand Index 0.38, fair agreement). Summer
heat is highest in coarse-grain industrial fabric. Morphology has modest but
genuine predictive skill ($R^2\approx0.09$), rising to $\approx0.27$ with coastal
and topographic context, which makes form secondary to the coastal gradient at
metropolitan scale. Combining heat with access, exposure and vulnerability
reverses a heat-only reading: the coolest fabric, the waterfront, becomes the
highest adaptation priority because it is the most vulnerable and exposed.
Weight robustness, parameters, layers and code are reported.
\end{abstract}

\noindent\textbf{Keywords:} morphometrics; explainable machine learning; SHAP;
Pareto frontier; TOPSIS; land-surface temperature; social vulnerability;
open-source GIS
